\documentclass[12pt,a4paper]{article}

\usepackage[a4paper,margin=2.2cm]{geometry}
\usepackage{fontspec}
\usepackage{polyglossia}
\usepackage{enumitem}
\usepackage{hyperref}

\setmainlanguage{russian}
\setotherlanguage{english}

\setmainfont[
  Path=fonts/,
  UprightFont=pt-astra-serif_regular.ttf,
  BoldFont=pt-astra-serif_bold.ttf,
  ItalicFont=pt-astra-serif_italic.ttf,
  BoldItalicFont=pt-astra-serif_bold-italic.ttf
]{PT Astra Serif}

\setmonofont[
  Path=fonts/fira-code/,
  UprightFont=FiraCode-Regular.ttf,
  BoldFont=FiraCode-Bold.ttf
]{Fira Code}

\hypersetup{
  hidelinks,
  pdfauthor={Нахатакян Артур Романович},
  pdftitle={Тезисы доклада для предзащиты ВКР}
}

\setlist[itemize]{leftmargin=1.4cm}
\setlist[enumerate]{leftmargin=1.4cm}
\setlength{\parindent}{0pt}
\setlength{\parskip}{0.55em}
\emergencystretch=2em

\begin{document}

\begin{center}
{\Large \textbf{Тезисы доклада для выступления на предзащите выпускной квалификационной работы}}\\[0.8em]
{\normalsize \textbf{Тема:} «Разработка системы управления контентом в сфере маркетинга на базе Strapi и Astro»}\\[0.4em]
{\normalsize \textbf{Обучающийся:} Нахатакян Артур Романович}\\
{\normalsize \textbf{Научный руководитель:} Государев Илья Борисович, к.п.н., доцент}
\end{center}

\section*{1. Актуальность темы}

Современный корпоративный маркетинговый сайт является не только представительской витриной,
но и рабочим инструментом бизнеса: через него публикуются страницы под рекламные кампании,
статьи, проекты, вакансии и формы захвата лидов. На практике такие сайты часто развиваются
быстрее, чем их контентная архитектура. В результате часть данных остается зашитой во
frontend-коде, публикация изменений требует участия разработчика, а локализация,
предварительный просмотр и \texttt{SEO}-настройки оказываются фрагментированными.

Актуальность работы обусловлена необходимостью построения \texttt{CMS-first} решения, в
котором контент, публикационный контур и публичная витрина рассматриваются как единая
инженерная система. Для этой задачи выбрана связка \texttt{Strapi 5} как ядра управления
контентом и \texttt{Astro 6} как производительной публичной витрины.

\section*{2. Объект, предмет, цель и задачи исследования}

\textbf{Объект исследования} --- процессы управления контентом корпоративного маркетингового
сайта.

\textbf{Предмет исследования} --- архитектурные и программные решения для построения
\texttt{CMS-first} веб-платформы, обеспечивающей управление контентом без постоянного
вмешательства в код публичной витрины.

\textbf{Цель работы} --- разработать систему управления контентом для корпоративного
маркетингового сайта на базе \texttt{Strapi} и \texttt{Astro}, обеспечивающую управление
страницами, статьями, проектами и вакансиями, поддержку локалей \texttt{ru/en},
предварительного просмотра, \texttt{SEO}-контур и воспроизводимый сценарий публикации.

Для достижения цели решались следующие задачи:
\begin{enumerate}
  \item проанализировать особенности маркетинговых сайтов и обосновать выбор
  \texttt{CMS-first} архитектуры;
  \item спроектировать модель данных и маршрутный контур системы;
  \item реализовать управление страницами и ключевыми контентными сущностями через
  \texttt{Strapi};
  \item реализовать публичную витрину на \texttt{Astro} с локализацией, \texttt{preview},
  \texttt{SEO} и \texttt{sitemap};
  \item организовать публикационный и эксплуатационный контур, включая
  \texttt{webhook -> rebuild};
  \item провести проверку ключевых пользовательских и редакторских сценариев.
\end{enumerate}

\section*{3. Архитектурная идея решения}

В работе реализован монорепозиторий на базе \texttt{Nx + pnpm}, включающий два основных
приложения:
\begin{itemize}
  \item \textbf{CMS-слой} на \texttt{Strapi 5}, отвечающий за хранение и редактирование
  контента;
  \item \textbf{frontend-слой} на \texttt{Astro 6}, отвечающий за публикацию контента на
  публичной витрине.
\end{itemize}

В \texttt{Strapi} сформирована модель данных для сущностей \texttt{global},
\texttt{home-page}, \texttt{page}, \texttt{article}, \texttt{author}, \texttt{project},
\texttt{vacancy}, \texttt{industry}, \texttt{job-role}, а также прикладных сущностей
\texttt{lead-submission} и \texttt{vacancy-application}. Для главной страницы и CMS-страниц
реализован конструктор контента на основе \texttt{Dynamic Zone}.

На стороне frontend реализован locale-prefixed маршрутный контур для storefront-core:
\texttt{/:locale/}, \texttt{/:locale/:slug/}, \texttt{/:locale/articles/...},
\texttt{/:locale/projects/...}. Это позволяет публиковать локализованный контент в обеих
локалях и централизованно управлять навигацией, \texttt{SEO} и страницами.

\section*{4. Реализованный функциональный результат}

На текущем этапе по коду и проверкам подтверждены следующие результаты:
\begin{itemize}
  \item реализовано управление главной страницей, CMS-страницами, статьями, проектами,
  вакансиями, маркетинговыми лидами и откликами на вакансии через \texttt{Strapi};
  \item реализован защищенный \texttt{preview mode} для \texttt{home-page}, \texttt{page},
  \texttt{article}, \texttt{project}, \texttt{vacancy};
  \item реализован \texttt{SEO/Open Graph}-контур для ключевых сущностей, а также генерация
  \texttt{sitemap};
  \item реализованы серверные маршруты для отправки маркетинговой lead form и формы отклика
  на вакансию с валидацией, \texttt{consent}, \texttt{honeypot} и ограничениями на загрузку
  файла;
  \item реализован эксплуатационный контур: \texttt{OpenAPI}-документация, versioned
  security model, \texttt{Docker}-конфигурация для CMS и managed
  \texttt{webhook -> rebuild} для обновления публичной витрины.
\end{itemize}

\section*{5. Проверка сценариев и количественные результаты}

Для проекта подготовлена воспроизводимая матрица приемки, разделяющая автоматизированные,
ручные и внешне ограниченные проверки.

По состоянию baseline от 22 мая 2026 года подтверждено:
\begin{itemize}
  \item публичные маршруты storefront-core, разделы статей, проектов и вакансий
  открываются корректно;
  \item \texttt{/api/preview} с неверным секретом возвращает \texttt{401}, а draft-preview
  корректно устанавливает preview-cookie и отдает страницы с \texttt{noindex};
  \item валидные отправки форм дают \texttt{201} и сохраняются в \texttt{SQLite};
  \item frontend-сборка проходит успешно и формирует \texttt{35} prerendered static routes;
  \item после сборки генерируются \texttt{sitemap-index.xml} и \texttt{sitemap-0.xml};
  \item в БД подтвержден активный \texttt{Frontend rebuild hook} на события
  \texttt{entry.publish} и \texttt{entry.unpublish};
  \item размер клиентского build output составляет около \texttt{2.3 МБ}.
\end{itemize}

Таким образом, результат работы подтвержден не только кодом, но и воспроизводимым
контуром проверок.

\section*{6. Практическая значимость}

Практическая значимость работы заключается в том, что разработанная система позволяет
перенести типовые редакторские операции из frontend-кода в \texttt{CMS}. Контент-менеджер
получает возможность управлять страницами, статьями, проектами и метаданными через
административный интерфейс, а разработчик перестает быть обязательным участником каждого
небольшого контентного изменения.

Для проекта в целом это дает:
\begin{itemize}
  \item централизованную модель данных;
  \item единый публикационный контур;
  \item поддержку локализации;
  \item управляемый \texttt{preview};
  \item подготовку публичной витрины к масштабируемому развитию.
\end{itemize}

\section*{7. Ограничения текущей версии и направления развития}

В ходе проверки зафиксированы и ограничения текущей версии системы.

Во-первых, архитектурный \texttt{ru/en} контур реализован, однако англоязычное наполнение
detail-разделов \texttt{articles/projects} пока зависит от фактического состояния dataset в
\texttt{CMS} и на baseline подтверждено не полностью.

Во-вторых, в текущем наборе \texttt{CMS SEO}-данных часть публичных страниц отдается с
\texttt{robots=noindex, nofollow}, что требует корректировки редакторского заполнения и
отдельной финальной проверки индексации.

В-третьих, browser-level аудит через \texttt{Lighthouse}/\texttt{axe} в данной сессии не
был завершен из-за ограничений внешней среды, поэтому accessibility и performance baseline
сейчас подтверждены преимущественно на структурном, build- и runtime-уровне.

В качестве направлений дальнейшего развития можно выделить:
\begin{itemize}
  \item расширение библиотеки \texttt{Dynamic Zone} блоков;
  \item доведение англоязычного контента до полной detail coverage;
  \item завершение browser-level аудита доступности и производительности;
  \item дальнейшее развитие редакционных ролей и workflow публикации.
\end{itemize}

\section*{8. Заключение}

В результате выполнения выпускной квалификационной работы разработан прототип
\texttt{CMS-first} системы управления контентом для корпоративного маркетингового сайта на
базе \texttt{Strapi} и \texttt{Astro}. Полученное решение охватывает управление контентом,
публичную витрину, локализацию, \texttt{preview}, \texttt{SEO}, формы взаимодействия с
пользователем и публикационный контур.

Практический результат работы состоит в том, что проект уже переведен от частично
кодо-центричного подхода к централизованной модели управления контентом. Это создает
инженерную основу для дальнейшего развития маркетинговой платформы как управляемой и
воспроизводимой системы.

\end{document}
